\section{The equivariant index calculation}\label{sec:index}

We complete the proof of Theorem~\ref{thm:division}.  Retain its notation,
and let \(r\) be the reduced rank of a finite right module over
\[
 \Delta=(\Delta _0\otimes_k F)\otimes_F\Delta _1.
\]
Our objective is to prove \(e_0e_1\mid r\).
Proposition~\ref{prop:rank-to-euler} supplies an \(H\)-equivariant coherent
class \(\beta\) on \(Z\times X\), and its reduction \(\beta'\) after a
complex embedding of \(k\).  It also supplies the two numerical
identities \eqref{eq:descent-fiber-index} and
\eqref{eq:descent-twist-divisibility}.  The first recovers \(r\) from an
Euler characteristic on \(Z\); the second gives divisibility of Euler
characteristics on \(Z\times X\).  We connect them by replacing low
dimensional tests in \(X\) with projective space.

\subsection{Equivariant maps into the smooth parameter space}

Put \(\mathbb T=(S^1)^M\).  The point characters defining the action on
the \(M\) blocks of the cone give a homomorphism
\(\rho:H\longrightarrow\mathbb T\); write \(H'=\rho(H)\).
The group \(\mathbb T\) need not preserve the cone, but its finite
subgroup \(H'\) does.  The \(H'\)-action and scalar multiplication by
\(S^1\) commute on the unit cone
\[
 Y_1=\{y\in Y^o:\|y\|=1\}.
\]
Here the norm is the standard Hermitian norm after the chosen complex
embedding.  Radial retraction from \(Y^o\) onto \(Y_1\) is equivariant
for both actions.  By Theorem~\ref{thm:parameters}, \(Y_1\) is
\((M+2h)\)-connected.

\begin{lemma}\label{lem:index-equivariant-map}
There is an \(H'\)-equivariant map
\[
 \Phi:\mathbb T\times\PP^h(\C)\longrightarrow X_{\sm},
\]
where \(H'\) acts by translation on \(\mathbb T\) and trivially on
\(\PP^h(\C)\).  It has an \(H'\times S^1\)-equivariant lift
\[
 \widetilde\Phi:\mathbb T\times S^{2h+1}\longrightarrow Y_1.
\]
If \(t_0\in\mathbb T\) is fixed and \(\Phi_0(z)=\Phi(t_0,z)\), then \(\Phi\)
is homotopic, as an ordinary map, to \(\Phi_0\) composed with projection
onto \(\PP^h(\C)\), and
\[
 \Phi_0^*\cO_X(-1)\simeq\cO_{\PP^h}(-1).
\]
There is also an \(H\)-equivariant map
\(\vartheta:P\longrightarrow\mathbb T\), for the action through \(\rho\).
\end{lemma}

\begin{proof}
The product of the principal bundles
\(\mathbb T\to\mathbb T/H'\) and
\(S^{2h+1}\to\PP^h(\C)\) is a principal \(H'\times S^1\)-bundle.
An equivariant lift \(\widetilde\Phi\) is a section of its associated
bundle with fiber \(Y_1\).  The base
\((\mathbb T/H')\times\PP^h(\C)\) has a finite CW structure of dimension
\(M+2h\).  Extend a section cell by cell: over an \(a\)-cell the
extension obstruction is a map \(S^{a-1}\to Y_1\), which is
null-homotopic by the connectivity just established.  The initial
choice over zero-cells and extension over one-cells are allowed
because \(Y_1\) is nonempty and path connected.  This constructs
\(\widetilde\Phi\), and passage to the scalar-circle quotient gives \(\Phi\).

To compare \(\Phi\) with \(\Phi_0\), forget \(H'\)-equivariance.
The lifts \(\widetilde\Phi(t,z)\) and \(\widetilde\Phi(t_0,z)\) are sections
of the same circle-associated bundle over
\(\mathbb T\times\PP^h(\C)\).  Construct a homotopy relative to its
two endpoints by extending over cells of the base times \([0,1]\).
This relative CW pair has dimension \(M+2h+1\), so the largest
obstruction belongs to \(\pi_{M+2h}(Y_1)\), which also vanishes.
Projectivizing this homotopy proves the assertion about \(\Phi\).
The circle-equivariant map
\(S^{2h+1}\to Y_1\) obtained at \(t_0\) identifies the associated
tautological line bundles, giving the stated isomorphism for \(\Phi_0\).

Finally, any character of \(H=G^{M-1}\) is a product of characters of
its factors.  Lemma~\ref{lem:picard-character-maps} realizes each
factor character by an equivariant circle-valued map on the
corresponding factor of \(P\).  Multiplying these maps realizes each
coordinate of \(\rho\).  Their product is \(\vartheta\).
\end{proof}

\subsection{An integral family of indices}

To bring \(\beta'\) to projective space, use the maps just constructed.
Writing \(z=(a,p)\in B\times P=Z\), define
\[
 \gamma:Z\times\PP^h(\C)\longrightarrow Z\times X_{\sm},
 \qquad
 \gamma(z,t)=(z,\Phi(\vartheta(p),t)).
\]
This map is \(H\)-equivariant, with trivial action on \(\PP^h\).
Since \(H\) acts freely on \(Z\), the quotient \(V=Z/H\) is smooth
projective, and \(\gamma\) induces a map
\[
 V\times\PP^h(\C)\longrightarrow (Z\times X_{\sm})/H.
\]
The ordinary homotopy from \(\Phi\) to \(\Phi_0\) will later compare
this pullback with an algebraic linear section.  For descent we use
\(\gamma\) itself, because \(\Phi_0\) need not be equivariant.
These maps are continuous, so the pullback class need not be algebraic.
We therefore need integration in topological \(K\)-theory that retains
integer coefficients, beyond the rational Chern-character formula.

We use \(K^0(T)\) for the Grothendieck group of complex topological
vector bundles on a compact space \(T\).  The group extends to the
noncompact smooth varieties below by the same vector-bundle
definition; a coherent class on such a variety defines a class in
this group by a finite locally free resolution.
On \(\PP^h(\C)\), put
\[
 u=1-[\cO_{\PP^h}(-1)]\in K^0(\PP^h(\C)).
\]
The projective bundle formula gives
\begin{equation}\label{eq:index-projective-k}
 K^0(\PP^h(\C))=\Z[u]/(u^{h+1}).
\end{equation}
In particular, the coefficients in this basis are integers, a fact
we need in addition to the rational Chern character.

The integrality statement below is a special case of the differentiable
Riemann--Roch theorem of Atiyah and Hirzebruch
\cite[Theorem~1 and Remark~(2)]{AtiyahHirzebruch59}.
We give the projective-space argument with an explicit choice of Thom
class.

\begin{lemma}\label{lem:index-integral-pushforward}
Let \(V\) be a smooth complex projective variety and let
\(\xi\in K^0(V\times\PP^h(\C))\).  There is a unique
\(I_V(\xi)\in K^0(\PP^h(\C))\) such that
\begin{equation}\label{eq:index-integral-pushforward}
 \ch(I_V(\xi))
   =\int_V\ch(\xi)\td(T_V).
\end{equation}
The integral on the right is integration over the \(V\)-factor.
\end{lemma}

\begin{proof}
Choose a projective embedding \(i:V\hookrightarrow\PP^a(\C)\),
with complex normal bundle \(\nu\).
The complex \(K\)-theory Thom class, followed by extension from a
tubular neighborhood, defines
\[
 i_!:K^0(V\times\PP^h(\C))
        \longrightarrow K^0(\PP^a(\C)\times\PP^h(\C)).
\]
We use the Koszul convention for this Thom class: its restriction
to the zero section of \(\nu\) is
\(\lambda_{-1}(\nu^\vee)\).
If \(x_1,\ldots,x_c\) are the Chern roots of \(\nu\), its Chern
character, divided by the ordinary cohomological Thom class, is
\[
 \prod_{\alpha=1}^c\frac{1-e^{-x_\alpha}}{x_\alpha}
       =\td(\nu)^{-1}.
\]
For completeness, this identity follows on the zero section from
\(\ch(\lambda_{-1}(\nu^\vee))=\prod_\alpha(1-e^{-x_\alpha})\)
and from the Euler class \(\prod_\alpha x_\alpha\).
The splitting principle verifies the identity universally for
sums of line bundles, where the product of the universal Chern
classes is a non-zero-divisor; the Thom isomorphism then gives
the asserted identity before restriction as well.
Consequently
\[
 \ch(i_!\xi)=i_*\bigl(\ch(\xi)\td(\nu)^{-1}\bigr).
\]
The exact tangent-normal sequence now implies
\begin{equation}\label{eq:index-ambient-integration}
 \int_{\PP^a}\ch(i_!\xi)\td(T_{\PP^a})
    =\int_V\ch(\xi)\td(T_V).
\end{equation}

The integral projective bundle formula identifies
\[
 K^0(\PP^a(\C)\times\PP^h(\C))
    =\Z[v,u]/(v^{a+1},u^{h+1}),\qquad
 v=1-[\cO_{\PP^a}(-1)].
\]
Write \(i_!\xi=\sum_{i=0}^a\sum_{j=0}^h c_{ij}v^iu^j\),
with \(c_{ij}\in\Z\).
The number
\[
 \int_{\PP^a}\ch(v^i)\td(T_{\PP^a})
    =\chi\bigl(\PP^a,(1-[\cO(-1)])^i\bigr)
\]
is an integer by Hirzebruch--Riemann--Roch, applied additively to
the finite sum of line-bundle classes \(v^i\).
Thus the left side of \eqref{eq:index-ambient-integration} is the
Chern character of an integral combination of \(1,u,\ldots,u^h\),
which gives \(I_V(\xi)\).
Uniqueness follows because the Chern character is injective on
the torsion-free group in \eqref{eq:index-projective-k}:
if \(H=c_1(\cO(1))\), then \(\ch(u)=1-e^{-H}\) has leading term \(H\),
so these \(h+1\) Chern characters are linearly independent over \(\Q\).
The projective bundle and Thom constructions used here are standard
complex \(K\)-theory, with the Koszul convention for the Thom class
specified above \cite[Proposition~2.24, Theorem~2.25,
and pp.~113--114]{Hatcher}; for the
Riemann--Roch formula see \cite[\S7]{BorelSerre58}.
\end{proof}

For a projective variety and a virtual coherent or algebraic
vector-bundle class, \(\chi\) will always mean the additive Euler
characteristic of that class.  We now apply the preceding lemma
to the equivariant class from Section~\ref{sec:descent}.

\begin{proposition}\label{prop:index-test}
In the setting of Theorem~\ref{thm:division}, let \(r\) be the
reduced rank of any finite right \(\Delta\)-module.
There is a smooth \(h\)-dimensional linear section
\(D_0\subset X_{\sm}\), a complete intersection of \(s\) hypersurfaces
of degree \(b\) in \(\PP^{h+s}(\C)\), and integers
\(q_0,\ldots,q_h\in e_1\Z\) such that \(q_0=\pm r\) and
\begin{equation}\label{eq:index-section-divisibility}
 \sum_{j=0}^h q_j\,
 \chi\!\left(D_0,
       (1-[\cO_{D_0}(-1)])^j\otimes\cO_{D_0}(l)\right)
       \in e_0e_1\Z
       \qquad(l\in\Z).
\end{equation}
\end{proposition}

\begin{proof}
Pull back the topological class of
\(\beta'|_{Z\times X_{\sm}}\) along \(\gamma\) and call the resulting
class \(\sigma\).
This pullback is permitted because the class is perfect on the
smooth open product.

The equivariant coherent class \(\beta'\) descends on the free smooth
quotient \((Z\times X_{\sm})/H\).  That quotient is smooth
quasi-projective, so the descended coherent class has a topological
vector-bundle class.  Pulling it back along the quotient of
\(\gamma\) gives
\(\tau\in K^0(V\times\PP^h(\C))\) whose pullback is \(\sigma\).
Since \(Z\to V\) is an étale covering of degree \(e_1\) and
\(T_Z\) is the pullback of \(T_V\), Lemma~\ref{lem:index-integral-pushforward}
gives
\begin{equation}\label{eq:index-family}
 \int_Z\ch(\sigma)\td(T_Z)
   =e_1\ch(I_V(\tau))
   =\ch\!\left(\sum_{j=0}^h q_j u^j\right),
 \qquad q_j\in e_1\Z.
\end{equation}

Nonequivariantly, Lemma~\ref{lem:index-equivariant-map} makes
\(\gamma\) homotopic to \((z,t)\mapsto(z,\Phi_0(t))\).
The equality \eqref{eq:index-family} therefore computes the same
fiber integral for this latter pullback of \(\beta'\).
Evaluate it at any point of \(\PP^h(\C)\).
Only \(q_0\) remains on the right, whereas
\eqref{eq:descent-fiber-index} gives \(\pm r\) on the left.
Thus \(q_0=\pm r\).

This already gives \(e_1\mid r\).  The stronger divisibility by
\(e_0e_1\) is encoded in the algebraic twist indices \(p(l)\) of
\eqref{eq:descent-twist-divisibility}.  A smooth linear section of \(X\)
retains that divisibility and can also be compared with the
projective-space test.  Put \(d=\dim X\).
Theorem~\ref{thm:parameters} gives
\(\operatorname{codim}_X(X\setminus X_{\sm})>h\).
A general linear section of codimension \(d-h\) therefore misses
the singular locus; Bertini gives a smooth section \(D_0\).
As \(X\) is a complete intersection of \(s\) degree-\(b\) equations,
\(D_0\) has the description in the statement.

Choose \(h+1\) linear forms with no common zero on \(D_0\);
these exist by general position since \(\dim D_0=h\).
They define a morphism
\[
 f:D_0\longrightarrow\PP^h(\C),
 \qquad f^*\cO_{\PP^h}(-1)\simeq\cO_{D_0}(-1).
\]
The inclusion \(j:D_0\hookrightarrow X_{\sm}\) and
\(\Phi_0\circ f\) have circle-equivariant lifts from the unit frame
bundle of this same line bundle to \(Y_1\).
Indeed, choose the indicated line-bundle isomorphisms and normalize
them by their positive fiber norms to identify the unit frame bundles.
The relative section-extension argument in
Lemma~\ref{lem:index-equivariant-map}, now over
\(D_0\times[0,1]\), makes those lifts homotopic:
the real dimension of this base is \(2h+1\), below the available
connectivity bound plus one.
Thus \(j\) and \(\Phi_0\circ f\) are homotopic.

Pulling \eqref{eq:index-family} back by \(f\) now identifies
the fiber integral on \(Z\times D_0\) with
\[
 \ch\!\left(
       \sum_{j=0}^h q_j(1-[\cO_{D_0}(-1)])^j
     \right).
 \tag{\(\ast\)}
\]
Define
\[
 p_{D_0}(l)=
 \chi\!\left(Z\times D_0,
       \beta'|_{Z\times D_0}\otimes\cO_{D_0}(l)\right),
\]
where restriction of a coherent class is derived.
Hirzebruch--Riemann--Roch on the smooth projective product, followed
by \((\ast)\), gives
\begin{equation}\label{eq:index-section-rr}
 p_{D_0}(l)=
 \sum_{j=0}^h q_j\,
 \chi\!\left(D_0,
       (1-[\cO_{D_0}(-1)])^j\otimes\cO_{D_0}(l)\right).
\end{equation}

It remains to retain the original divisibility when passing to
this section.  Recall from \eqref{eq:descent-twist-divisibility}
the integers
\[
 p(l)=\chi(Z\times X,\beta\otimes\cO_X(l))/e_0
       \in e_0e_1\Z.
\]
The \(d-h\) hyperplanes defining \(D_0\) have a Koszul resolution
on \(X\): along \(D_0\) they cut transversely in the smooth locus,
and off their common zero locus the complex is exact because
one of the equations is locally a unit.
Tensoring this resolution derivedly with \(\beta\), and then
passing to the complex embedding and the matrix-idempotent
reduction, yields
\begin{equation}\label{eq:index-koszul-difference}
 p_{D_0}(l)
   =\sum_{a=0}^{d-h}(-1)^a\binom{d-h}{a}p(l-a)
   \in e_0e_1\Z.
\end{equation}
Here cohomology commutes with the field extension, and the
matrix-idempotent reduction divides every cohomology dimension by
\(e_0\).  The particular hyperplanes need only be defined after
that extension: the terms on the right of
\eqref{eq:index-koszul-difference} are the original integer Euler
characteristics.  Combining
\eqref{eq:index-section-rr} and \eqref{eq:index-koszul-difference}
proves \eqref{eq:index-section-divisibility}.
\end{proof}

For the integers in Proposition~\ref{prop:index-test}, put
\[
 Q_*(u)=\sum_{j=0}^h(q_j/e_1)u^j\in\Z[u].
\]
Its constant coefficient is \(Q_*(0)=\pm r/e_1\), and all its twisted
Euler characteristics on \(D_0\) are divisible by \(e_0\).
The remaining step extracts the same divisibility for \(Q_*(0)\).
The dimension \(h\) was chosen precisely to allow the following
finite truncation argument.

\subsection{Truncated polynomial congruences}

\begin{lemma}\label{lem:index-padic-truncation}
Let \(p\) be a prime, let \(v\ge1\), and let
\(A,Q\in\Z[u]\).  Suppose the reduction of \(A\) modulo \(p\)
has exact order \(\nu\ge0\) at \(u=0\).
If \(h\ge v\nu\) and
\[
 A(u)Q(u)=0\pmod{(p^v,u^{h+1})},
\]
then \(p^v\mid Q(0)\).
\end{lemma}

\begin{proof}
We successively construct integer polynomials \(Q_i\),
for \(0\le i\le v\), with \(Q_0=Q\), such that
\[
 Q_i(0)=Q(0)/p^i,\qquad
 A Q_i=0\pmod{(p^{v-i},u^{h-i\nu+1})}
       \quad(0\le i<v).
\]
Suppose \(i<v\) and \(Q_i\) has been constructed.
Modulo \(p\), write \(A=u^\nu B\), where \(B(0)\ne0\).
The power series \(B\) is invertible over \(\Z/p\), so the displayed
congruence forces every coefficient of \(Q_i\) through degree
\[
 m_i=h-(i+1)\nu
\]
to be divisible by \(p\).  This degree is nonnegative because
\(h\ge v\nu\).
Let \(Q_{i+1}\) be the truncation of \(Q_i\) through degree \(m_i\),
divided by \(p\).
It is an integer polynomial and has the required constant
coefficient.  Since
\[
 Q_i-pQ_{i+1}\in u^{m_i+1}\Z[u],
\]
restricting the old congruence to degrees at most \(m_i\) gives
\[
 p A Q_{i+1}=0\pmod{(p^{v-i},u^{m_i+1})}.
\]
Coefficientwise division by \(p\) gives the next congruence
modulo \(p^{v-i-1}\).
At \(i=v-1\), the construction shows directly that
\(Q(0)/p^{v-1}\) is divisible by \(p\).
This proves the conclusion.  The same induction includes
\(\nu=0\), when no degrees are lost.
\end{proof}

\begin{proof}[Completion of the proof of Theorem~\ref{thm:division}]
Use \(D_0\), \(q_0,\ldots,q_h\), and \(Q_*\) just constructed for
the chosen module.
If \(e_0=1\), then \(q_0=\pm r\in e_1\Z\) already gives the required
divisibility, so assume \(e_0>1\).
The Hilbert series of the coordinate ring of the complete
intersection \(D_0\) is
\[
 \frac{(1-t^b)^s}{(1-t)^{h+s+1}}.
\]
Euler characteristics of sufficiently positive twists agree
with Hilbert functions.  Moreover, multiplication by
\((1-[\cO(-1)])^j\) gives the same finite difference as
multiplication of generating functions by \((1-t)^j\), apart
from finitely many initial coefficients.  The coefficients in
these differences are integers.  It follows from
\eqref{eq:index-section-rr} that
\begin{equation}\label{eq:index-hilbert-series}
 \sum_{l\ge0}\frac{p_{D_0}(l)}{e_1}t^l
   =Q_*(1-t)\frac{(1-t^b)^s}{(1-t)^{h+s+1}}+R(t),
 \qquad R(t)\in\Z[t].
\end{equation}
By \eqref{eq:index-koszul-difference}, every coefficient on the
left is divisible by \(e_0\).

Multiply \eqref{eq:index-hilbert-series} by
\((1-t)^{h+s+1}\) and reduce coefficients modulo \(e_0\).
The resulting right side is a polynomial that vanishes
coefficientwise.  We may therefore substitute \(t=1-u\), obtaining
\[
 Q_*(u)\bigl(1-(1-u)^b\bigr)^s
     =0\pmod{(e_0,u^{h+s+1})}.
\]
Define
\[
 A_b(u)=\frac{1-(1-u)^b}{u}\in\Z[u].
\]
Cancelling the common monomial \(u^s\) yields
\begin{equation}\label{eq:index-final-congruence}
 A_b(u)^sQ_*(u)=0\pmod{(e_0,u^{h+1})}.
\end{equation}
This is coefficientwise cancellation of a monomial, valid over
\(\Z/e_0\) even when \(e_0\) is composite.  When \(s=0\), the
factor is \(1\), and the constant coefficient of
\eqref{eq:index-final-congruence} immediately gives
\(e_0\mid Q_*(0)\).

For the general case, fix a prime \(p\mid e_0\) and write
\(v=v_p(e_0)\), \(b=p^a c\), where \(p\nmid c\).
Over \(\Z/p\),
\[
 1-(1-u)^b
   =1-(1-u^{p^a})^c
   =c\,u^{p^a}+\text{terms of higher degree}.
\]
Thus \(A_b(u)^s\) has exact order
\(\nu=s(p^a-1)\) modulo \(p\), with nonzero leading coefficient.
The choice of \(h\) in Theorem~\ref{thm:division} gives
\(h\ge v\nu\).
Apply Lemma~\ref{lem:index-padic-truncation} to
\eqref{eq:index-final-congruence}, reduced modulo \(p^v\), to obtain
\(p^v\mid Q_*(0)\).
This also covers primes not dividing \(b\), for which \(\nu=0\);
no assumption that \(b\) is a prime power has been made.
Doing this for every prime divisor of \(e_0\) gives
\[
 e_0\mid Q_*(0),\qquad
 e_0e_1\mid q_0=\pm r.
\]
The module was arbitrary.  A nonzero proper right ideal of the
central simple algebra \(\Delta\), whose degree is \(e_0e_1\),
would have reduced rank strictly between \(0\) and \(e_0e_1\).
The divisibility just proved excludes such an ideal, so
\(\Delta\) is division.
\end{proof}
